% Generated draft: requires mathematical and source review.
\documentclass[11pt]{article}
\usepackage[margin=25mm]{geometry}
\usepackage{fontspec}
\setmainfont{Latin Modern Roman}
\usepackage{amsmath,amssymb,mathtools,unicode-math}
\setmathfont{Latin Modern Math}
\usepackage{enumitem,needspace,xurl,hyperref}
\setlength{\parindent}{0pt}
\setlength{\parskip}{5pt}
\setlength{\emergencystretch}{3em}
\newcommand{\N}{\mathbb{N}}
\newcommand{\Z}{\mathbb{Z}}
\newcommand{\Q}{\mathbb{Q}}
\newcommand{\R}{\mathbb{R}}
\newcommand{\C}{\mathbb{C}}
\newcommand{\F}{\mathbb{F}}
\newcommand{\A}{\mathbb{A}}
\newcommand{\PP}{\mathbb{P}}
\newcommand{\E}{\mathbb{E}}
\newcommand{\eps}{\varepsilon}
\newcommand{\dd}{\,\mathrm{d}}
\begin{document}
\section*{UnsolvedMath ID: 800896}
\textbf{Source problem code:} catalogue entry
\subsection{Definitions and mathematical statement}
The supplied record does not define \(g_{\Omega}\), specify the ambient complex dimension, or explicitly locate the pole \(w\). The following formulation uses the conventional interpretation as the pluricomplex Green function with an interior logarithmic pole; these interpretations require confirmation against the source.

Let \(n\geq 1\), let \(\Omega\subset\C^n\) be a bounded domain with \(C^{\infty}\) boundary, and let \(w\in\Omega\). Strict pseudoconvexity means that there is a smooth real defining function \(\rho\) near \(\partial\Omega\), with \(\Omega=\{\rho<0\}\) locally and \(\dd\rho\ne0\) on the boundary, such that
\[
\sum_{j,k=1}^{n}\frac{\partial^2\rho}{\partial z_j\partial\overline z_k}(p)\,\xi_j\overline{\xi_k}>0
\]
for every \(p\in\partial\Omega\) and every nonzero \(\xi\in\C^n\) satisfying \(\sum_j\rho_{z_j}(p)\xi_j=0\). A function is plurisubharmonic if it is upper semicontinuous and its restriction to every complex line is subharmonic wherever defined. Write \(\operatorname{PSH}(\Omega)\) for these functions, and define
\[
g_{\Omega}(z,w)=\left(\sup\left\{u(z):u\in\operatorname{PSH}(\Omega),\ u\leq0,\ u(\zeta)\leq\log\|\zeta-w\|+O(1)\text{ as }\zeta\to w\right\}\right)^*.
\]
Here \(O(1)\) denotes a quantity bounded near \(w\), and the star denotes upper semicontinuous regularization.

Under this interpretation, the question is whether there exist \(n\), \(\Omega\), and \(w\) as above such that
\[
g_{\Omega}(\cdot,w)\notin C^2(\Omega\setminus\{w\}).
\]
The class \(C^2\) refers to continuous real partial derivatives through order two on the punctured domain, with no boundary regularity requirement. The supplied statement reports the known regularity \(g_{\Omega}(\cdot,w)\in C^{1,1}(\Omega\setminus\{w\})\); if this notation is intended locally, it means continuous first real derivatives that are locally Lipschitz. This regularity assertion is retained as a supplied premise, not independently verified.

The source does not specify whether the dimension is fixed or allowed to vary, whether ``smoothly bounded'' explicitly includes boundedness, or whether its \(C^{1,1}\) notation imposes global estimates. The bounded, smooth-boundary interpretation and the existential quantifier over dimension used above are provisional.
\subsection{Short English statement}
Under the conventional pluricomplex Green function interpretation, ask whether some smoothly bounded strictly pseudoconvex domain has a Green function that is not twice continuously differentiable away from its interior pole. The supplied record reports \(C^{1,1}\) regularity and labels the question open; its notation, dimension, and domain conventions require source confirmation.
\subsection{Sources}
\begin{itemize}
\item Local UnsolvedMath registry, permanent ID \(800896\), exact source problem code \texttt{catalogue entry}; supplied complete source record in the AMR Open Problem Lists collection. The registry labels the problem \texttt{open}; this is a supplied status, not an independently verified literature assessment.
\item Dinew, Guedj, Zeriahi, \emph{Open Problems in Pluripotential Theory} (2015), Question 27, \url{https://arxiv.org/abs/1511.00705}. The supplied record identifies the extraction method as \texttt{arxiv\_latex} and records access on 2026-10-16; this access information is supplied provenance, not a claim of independent consultation.
\item The title and background attribute the question or its context to Blocki. The registry lists Dinew, Guedj, Zeriahi as proposers and 2015 as the proposed year. No additional literature references or literature-verification date are supplied.
\end{itemize}
\end{document}
